\documentclass[12pt, a4paper]{article}

% ── Packages ──────────────────────────────────────────────────────────────────
\usepackage[margin=2.5cm]{geometry}
\usepackage{amsmath, amssymb, amsthm}
\usepackage{booktabs}
\usepackage{array}
\usepackage{longtable}
\usepackage{hyperref}
\usepackage{xcolor}
\usepackage{titlesec}
\usepackage{fancyhdr}
\usepackage{graphicx}
\usepackage{float}
\usepackage{enumitem}
\usepackage{mathtools}
\usepackage{bbm}
\usepackage{caption}
\usepackage{subcaption}
\usepackage[round]{natbib}

% ── Colours ───────────────────────────────────────────────────────────────────
\definecolor{darkblue}{RGB}{0,51,102}
\definecolor{midblue}{RGB}{0,102,179}
\definecolor{lightgray}{RGB}{245,245,245}

% ── Hyperref setup ────────────────────────────────────────────────────────────
\hypersetup{
    colorlinks=true,
    linkcolor=darkblue,
    urlcolor=midblue,
    citecolor=darkblue,
    pdfauthor={Risk Management},
    pdftitle={Margin Requirement Methodology}
}

% ── Header / Footer ───────────────────────────────────────────────────────────
\pagestyle{fancy}
\fancyhf{}
\rhead{\textcolor{darkblue}{\small Margin Requirement Methodology v1.0}}
\lhead{\textcolor{darkblue}{\small Confidential}}
\cfoot{\thepage}

% ── Section formatting ────────────────────────────────────────────────────────
\titleformat{\section}{\large\bfseries\color{darkblue}}{{\thesection}}{1em}{}[\titlerule]
\titleformat{\subsection}{\normalsize\bfseries\color{midblue}}{\thesubsection}{1em}{}

% ── Theorem environments ──────────────────────────────────────────────────────
\newtheorem{definition}{Definition}[section]
\newtheorem{proposition}{Proposition}[section]
\newtheorem{remark}{Remark}[section]

% ── Custom commands ───────────────────────────────────────────────────────────
\newcommand{\MR}{\mathrm{MR}}
\newcommand{\IM}{\mathrm{IM}}
\newcommand{\MM}{\mathrm{MM}}
\newcommand{\E}{\mathbb{E}}
\newcommand{\Prob}{\mathbb{P}}
\newcommand{\R}{\mathbb{R}}
\newcommand{\VaR}{\mathrm{VaR}}
\newcommand{\CVaR}{\mathrm{CVaR}}
\newcommand{\LTV}{\mathrm{LTV}}
\DeclareMathOperator*{\argmin}{arg\,min}

% ══════════════════════════════════════════════════════════════════════════════
\begin{document}

% ── Title Page ────────────────────────────────────────────────────────────────
\begin{titlepage}
    \centering
    \vspace*{2cm}
    {\color{darkblue}\rule{\linewidth}{1.5pt}}\\[0.4cm]
    {\Huge\bfseries\color{darkblue} Margin Requirement Methodology\\[0.3cm]
    for Margin Trading}\\[0.4cm]
    {\color{darkblue}\rule{\linewidth}{1.5pt}}\\[2cm]

    {\large\textbf{Version:} 1.0}\\[0.3cm]
    {\large\textbf{Classification:} Confidential}\\[0.3cm]
    {\large\textbf{Owner:} Risk Management / Quantitative Risk}\\[0.3cm]
    {\large\textbf{Date:} \today}\\[3cm]

    \begin{abstract}
    \noindent
    This document establishes the two-layer margin requirement framework
    underpinning the margin trading product. The framework is structured as
    a bottom-up hierarchy: a \emph{Layer~1} instrument-by-instrument floor
    grounded in the FINRA/Regulation~T paradigm, and a \emph{Layer~2}
    portfolio-level overlay that accounts for cross-asset correlations,
    diversification effects, concentration risk, and dynamic regime adjustment.
    The instrument floor provides an unconditional, rule-based minimum that
    is always enforced; the portfolio overlay may only tighten margin
    requirements relative to that floor, never relax them below it.
    This document serves as the methodological foundation for the
    Loan-to-Value~(LTV) Methodology (\emph{see} \texttt{trb\_method.tex}),
    which translates the margin requirement framework into operational LTV
    thresholds, haircut schedules, and liquidation triggers.
    \end{abstract}

    \vfill
    {\small\color{gray} This document is confidential and intended solely for use
    by authorised personnel of the Platform and the Lender. Reproduction or
    distribution without written consent is prohibited.}
\end{titlepage}

% ── Table of Contents ─────────────────────────────────────────────────────────
\tableofcontents
\newpage

% ══════════════════════════════════════════════════════════════════════════════
\section{Purpose and Scope}
\label{sec:scope}

\subsection{Purpose}
This document establishes the quantitative framework for determining margin
requirements across the margin trading product. It defines the methodology
for both the instrument-level and portfolio-level treatment of margin,
providing a single authoritative reference for:
\begin{itemize}[leftmargin=1.5em]
    \item the regulatory and conceptual basis for instrument-by-instrument
    margin requirements (Layer~1);
    \item the derivation of portfolio-level margin adjustments incorporating
    correlation, diversification, and concentration effects (Layer~2);
    \item the interaction rule between the two layers, ensuring the
    portfolio overlay can only tighten the instrument floor;
    \item the governance, recalibration, and validation requirements for
    each layer.
\end{itemize}

\noindent The LTV Methodology document (\texttt{trb\_method.tex}) references
this framework directly and translates its outputs into operational LTV
thresholds, haircut schedules, dynamic regime adjustments, and liquidation
triggers.

\subsection{Scope}
This methodology applies to all margined positions held as collateral in
segregated custody accounts for retail and professional clients. It covers
equities, broad-market ETFs, and fixed income instruments that meet
collateral eligibility criteria.

\noindent The following are \emph{explicitly out of scope}:
\begin{itemize}[leftmargin=1.5em]
    \item \textbf{Derivatives margining} -- futures, options, swaps, and any
    other exchange-traded or OTC derivative instruments. Clients are restricted
    to non-derivative instruments only; accordingly, concepts specific to
    derivatives margining (variation margin in the EMIR sense, premium margin,
    futures mark-to-market settlement) do not apply and are not addressed.
    \item \textbf{Premium margin} -- the margin component covering the current
    market value of short option positions. Entirely irrelevant given the
    non-derivative client universe.
    \item \textbf{Short-selling margin requirements.}
    \item \textbf{Unsecured lending.}
\end{itemize}

\subsection{Regulatory Context}
The framework draws on the following regulatory standards:
\begin{itemize}[leftmargin=1.5em]
    \item \textbf{Regulation T (Reg~T)} -- Federal Reserve regulation
    governing the extension of credit by brokers and dealers; establishes
    a 50\% initial margin requirement for most equity securities and forms
    the basis of the instrument-level floor in Layer~1.
    \item \textbf{FINRA Rule 4210} -- FINRA maintenance margin requirements
    (minimum 25\% equity for long positions); provides the maintenance
    floor parameters for Layer~1.
    \item \textbf{MiFID~II / MiFIR} -- suitability and leverage requirements
    for retail and professional clients in the EU.
    \item \textbf{CRR~II / Basel~III} -- Financial Collateral Comprehensive
    Method (FCCM) for regulatory haircut floors, referenced in Layer~1
    haircut calibration.
    \item \textbf{ESMA Guidelines on Leverage} -- retail client leverage
    caps by asset class, acting as binding ceilings on both layers.
\end{itemize}

% ══════════════════════════════════════════════════════════════════════════════
\section{Key Definitions and Margin Concepts}
\label{sec:definitions}

\begin{definition}[Current Liquidating Margin]
The \emph{current liquidating margin} (CLM) is the net equity value of a
client account if all open positions were closed at current market prices.
It is defined as:
\[
\mathrm{CLM}_t = V_t^{\mathrm{adj}} - L_t,
\]
where $V_t^{\mathrm{adj}}$ is the adjusted collateral value and $L_t$ is
the outstanding loan balance including accrued interest. CLM is the primary
measure of a client's cushion against a margin call and is monitored
continuously. It is the non-derivative equivalent of the equity balance
concept in exchange-margined accounts.
\end{definition}

\begin{definition}[Mark-to-Market Collateral Revaluation]
The daily (and intraday) revaluation of collateral positions at current
market prices. Changes in $V_t^{\mathrm{adj}}$ directly affect CLM and
may trigger margin calls when CLM falls below the maintenance margin
requirement. This is the functional analogue of \emph{variation margin} as
used in derivatives markets (EMIR), adapted to the non-derivative collateral
context. In this framework it is governed by the LTV calculation schedule
defined in \texttt{trb\_method.tex} Section~4.3, which specifies real-time,
intraday, and end-of-day revaluation frequencies and price source hierarchies.
\emph{Variation margin in the EMIR/futures sense -- daily cash settlement of
mark-to-market P\&L on derivative positions -- does not apply here, as
clients trade non-derivative instruments only.}
\end{definition}

\begin{definition}[Initial Margin]
The minimum equity a client must deposit at the time of loan origination,
expressed as a fraction of the position market value (see
Section~\ref{subsec:initial_margin}).
\end{definition}

\begin{definition}[Maintenance Margin]
The minimum equity a client must maintain at all times. A breach triggers
a margin call. Defined at the instrument level in Layer~1 and at the
portfolio level in Layer~2 (see Sections~\ref{subsec:maintenance_margin}
and \ref{sec:layer2}).
\end{definition}

% ══════════════════════════════════════════════════════════════════════════════
\section{Framework Overview}
\label{sec:overview}

The margin requirement $\MR_t$ for a client portfolio at time $t$ is
determined as the maximum of the two layer outputs:
\begin{equation}
    \MR_t = \max\!\left(\MR_t^{(1)},\; \MR_t^{(2)}\right),
    \label{eq:mr_total}
\end{equation}
where $\MR_t^{(1)}$ is the Layer~1 instrument-level aggregate and
$\MR_t^{(2)}$ is the Layer~2 portfolio-level requirement. This construction
guarantees that the portfolio overlay can \emph{only increase} the total
margin requirement relative to the instrument floor -- it can never grant
a credit that reduces it below the regulatory minimum.

\begin{figure}[H]
\centering
\begin{tabular}{|c|}
\hline
\\
\textbf{Layer 2: Portfolio Overlay} \\
{\small Correlation $\cdot$ Diversification $\cdot$ Concentration $\cdot$ Regime} \\
\\
\hline
\\
\textbf{Layer 1: Instrument Floor} \\
{\small Reg T / FINRA~4210 $\cdot$ Asset-Class Haircuts $\cdot$ FCCM Floors} \\
\\
\hline
\end{tabular}
\caption{Two-layer margin requirement structure. Layer~1 provides the
unconditional floor; Layer~2 may only add to it.}
\label{fig:layers}
\end{figure}

% ══════════════════════════════════════════════════════════════════════════════
\section{Layer 1: Instrument-by-Instrument Margin (Reg~T Floor)}
\label{sec:layer1}

\subsection{Conceptual Basis}
Layer~1 follows the Regulation~T and FINRA Rule~4210 paradigm: each
position is assessed independently against prescribed initial and maintenance
margin rates. No netting or correlation benefit is recognised at this stage.
The aggregate Layer~1 requirement is the sum of individual instrument
requirements:
\begin{equation}
    \MR_t^{(1)} = \sum_{i=1}^{N} \MR_{i,t}^{(1)}.
    \label{eq:layer1_aggregate}
\end{equation}

\subsection{Initial Margin Requirement}
\label{subsec:initial_margin}

The initial margin requirement establishes the minimum client equity that
must be in place \emph{at the time a margin loan is originated or a new
position is added to the collateral pool}. It acts as the first line of
defence against credit loss: by requiring the client to fund a minimum
fraction of the position from their own equity, the Platform limits its
maximum loan exposure relative to collateral value from the outset.

The initial margin requirement for a long position in instrument $i$ is:
\begin{equation}
    \IM_{i,t} = m_i^{\IM} \cdot P_{i,t} \cdot Q_i,
    \label{eq:initial_margin}
\end{equation}
where $m_i^{\IM} \in (0,1]$ is the initial margin rate applicable to
instrument $i$ and $P_{i,t} \cdot Q_i$ is the mark-to-market value of
the position at origination.

\subsubsection*{Regulatory Floor}
Under Regulation~T, $m_i^{\IM} \geq 0.50$ for marginable equity securities,
meaning the client must fund at least 50\% of the purchase price from equity.
The Platform applies house initial margin rates at or above this floor,
calibrated by asset class as detailed in Table~\ref{tab:margin_rates}.
For fixed income instruments, the Reg~T floor does not prescribe a specific
rate; the Platform applies rates consistent with CRR~II FCCM supervisory
haircuts scaled to an equivalent margin basis.

\subsubsection*{Relationship to LTV at Origination}
The initial margin rate determines the maximum LTV at which a new loan may
be originated. Specifically, the maximum initial LTV is:
\begin{equation}
    \LTV^{\mathrm{init}}_i = 1 - m_i^{\IM},
    \label{eq:ltv_init}
\end{equation}
so a 50\% initial margin rate implies a maximum origination LTV of 50\%.
For a multi-asset portfolio, the effective initial LTV cap is the
value-weighted average of per-instrument caps, subject to the portfolio-level
regime adjustment defined in Layer~2 (Section~\ref{sec:layer2}) and the
regime-dependent initial LTV schedule in \texttt{trb\_method.tex} Section~5.4.

\subsubsection*{Recalculation Triggers}
Although the initial margin requirement is primarily assessed at origination,
it is recalculated and reapplied in the following circumstances:
\begin{itemize}[leftmargin=1.5em]
    \item A client adds a new position or increases an existing position.
    \item A position is reclassified to a higher-risk asset class (e.g.\
    following a credit rating downgrade or a reduction in ADTV below a
    liquidity threshold).
    \item A regime transition to Stress is declared, triggering the lower
    initial LTV cap for newly originated or materially increased positions
    (existing positions are not force-cured solely on regime change).
\end{itemize}

\subsection{Maintenance Margin Requirement}
\label{subsec:maintenance_margin}

The maintenance margin requirement is the minimum equity the client must
hold at all times:
\begin{equation}
    \MM_{i,t} = m_i^{\MM} \cdot P_{i,t} \cdot Q_i,
    \label{eq:maintenance_margin}
\end{equation}
where $m_i^{\MM}$ is the maintenance margin rate. Under FINRA Rule~4210,
$m_i^{\MM} \geq 0.25$ for long equity positions. The Platform applies
enhanced house maintenance rates as detailed in Table~\ref{tab:margin_rates}.

\subsection{Asset-Class Margin Rate Schedule}
\label{subsec:margin_rates}

\begin{table}[H]
\centering
\caption{Instrument-Level Margin Rate Schedule (Layer~1)}
\label{tab:margin_rates}
\begin{tabular}{llcc}
\toprule
\textbf{Asset Class} & \textbf{Subcategory}
    & \textbf{Initial Margin $m^{\IM}$}
    & \textbf{Maintenance Margin $m^{\MM}$} \\
\midrule
Government Bonds & AAA--AA, maturity $< 1$Y     & 2\%   & 1\%   \\
                 & AAA--AA, 1--5Y                & 5\%   & 3\%   \\
                 & AAA--AA, 5--10Y               & 8\%   & 5\%   \\
                 & A--BBB, any maturity          & 15\%  & 10\%  \\
IG Corporate Bonds & 0--5Y                       & 15\%  & 10\%  \\
                   & 5Y+                         & 20\%  & 15\%  \\
Large-Cap Equities & MSCI World, high liquidity  & 50\%  & 30\%  \\
Mid-Cap Equities   & Mid liquidity               & 60\%  & 35\%  \\
Small-Cap Equities & Low liquidity               & 75\%  & 50\%  \\
ETFs (broad)       & Major index replication     & 50\%  & 30\%  \\
ETFs (thematic)    & Sector / factor exposure    & 60\%  & 35\%  \\
\bottomrule
\end{tabular}
\end{table}

\noindent Initial margin rates at or above 50\% for equities are consistent
with Reg~T. Rates for fixed income reflect FCCM supervisory haircut floors
scaled to an equivalent margin basis.

\subsection{Layer 1 Margin Call Trigger}

A Layer~1 margin call is triggered at the instrument level when the equity
in a single-position account falls below the maintenance margin requirement,
i.e.\ when:
\[
E_{i,t} = P_{i,t}\,Q_i - D_{i,t} < \MM_{i,t},
\]
where $D_{i,t}$ is the debit balance (amount borrowed) attributed to
position $i$ on a pro-rata basis. The required deposit to restore the
position to the initial margin level is:
\[
\Delta_i^{\mathrm{call}} = \IM_{i,t} - E_{i,t}.
\]

\subsection{FX and Non-EUR Positions}
\label{subsec:fx_layer1}

For assets denominated in currency $c \neq \mathrm{EUR}$, the margin
requirement is computed on the EUR-equivalent market value, and an
additional FX margin add-on is applied:
\begin{equation}
    \MR_{i,t}^{(1),\mathrm{FX}}
    = m_i^{\IM} \cdot P_{i,t}^{\mathrm{EUR}} \cdot Q_i
    \;+\; m_{\mathrm{FX}(c)} \cdot P_{i,t}^{\mathrm{EUR}} \cdot Q_i,
    \label{eq:layer1_fx}
\end{equation}
where $m_{\mathrm{FX}(c)}$ is the FX margin rate for currency pair
$c$/EUR, calibrated consistently with the haircut formula in
\texttt{trb\_method.tex} Section~3.2.

% ══════════════════════════════════════════════════════════════════════════════
\section{Layer 2: Portfolio-Level Margin Overlay}
\label{sec:layer2}

\subsection{Conceptual Basis}

Layer~2 assesses margin requirements at the portfolio level, recognising
that the true risk of a collateral portfolio is not merely the sum of
individual position risks. Two distinct effects are captured:

\begin{enumerate}[leftmargin=1.5em]
    \item \textbf{Diversification benefit:} Low or negative pairwise
    correlations reduce portfolio-level risk below the sum-of-parts.
    Layer~2 can recognise this benefit, subject to a floor equal to the
    Layer~1 aggregate (equation~\eqref{eq:mr_total}).
    \item \textbf{Concentration and correlation penalty:} Highly correlated
    or concentrated portfolios may exhibit risk \emph{above} the sum-of-parts
    during stress. Layer~2 applies a multiplicative penalty in such cases,
    increasing margin requirements beyond the Layer~1 floor.
\end{enumerate}

\noindent The Layer~2 margin requirement is defined as:
\begin{equation}
    \MR_t^{(2)} = \VaR_\alpha\!\left(\Delta V_t;\,\hat{\Sigma}_t,\,\Delta t\right)
    \cdot \Phi(\bar{\rho}_t, \kappa_t),
    \label{eq:layer2}
\end{equation}
where $\VaR_\alpha(\cdot)$ is the portfolio VaR at confidence level $\alpha$,
estimated under the covariance structure $\hat{\Sigma}_t$, and
$\Phi(\bar{\rho}_t, \kappa_t)$ is a concentration and correlation multiplier
defined in Section~\ref{subsec:concentration_multiplier}.

\subsection{Portfolio VaR Estimation}
\label{subsec:portfolio_var}

The portfolio profit-and-loss over the liquidation horizon $\Delta t$ is
approximated as:
\[
\Delta V_t \approx \sum_{i=1}^{N} w_i \cdot P_{i,t} \cdot Q_i \cdot r_{i,t+\Delta t},
\]
where $w_i = (1-h_i)$ is the haircut-adjusted weight of instrument $i$
and $r_{i,t+\Delta t}$ is the return over $[t, t+\Delta t]$.

\noindent Returns are modelled as jointly normal with zero mean and
covariance matrix $\hat{\Sigma}_t \cdot \Delta t$, where $\hat{\Sigma}_t$
is estimated via EWMA with decay factor $\lambda = 0.94$ on a 252-day
window. The portfolio variance is:
\begin{equation}
    \hat{\sigma}_{p,t}^2 = \mathbf{w}^\top \hat{\Sigma}_t\, \mathbf{w},
    \label{eq:port_var}
\end{equation}
where $\mathbf{w} = (w_1 P_{1,t} Q_1, \ldots, w_N P_{N,t} Q_N)^\top /
V_t^{\mathrm{adj}}$ is the vector of adjusted value weights.

\noindent The parametric portfolio VaR at confidence level $\alpha$
over horizon $\Delta t$ is:
\begin{equation}
    \VaR_\alpha\!\left(\Delta V_t\right)
    = z_\alpha\,\hat{\sigma}_{p,t}\,\sqrt{\Delta t}\,V_t^{\mathrm{adj}},
    \label{eq:port_var_formula}
\end{equation}
where $z_\alpha = \Phi^{-1}(\alpha)$ is the standard normal quantile
(e.g.\ $z_{0.99} = 2.326$ for $\alpha = 99\%$).

\begin{remark}
    For portfolios with significant fat-tail exposures, the parametric
    normal VaR in equation~\eqref{eq:port_var_formula} may underestimate
    tail risk. In such cases the portfolio VaR is estimated via Monte Carlo
    simulation using a multivariate $t$-distribution with fitted degrees of
    freedom $\nu$, as specified in the stress testing section of
    \texttt{trb\_method.tex}.
\end{remark}

\subsection{Concentration and Correlation Multiplier}
\label{subsec:concentration_multiplier}

The multiplier $\Phi$ is composed of two additive adjustments:
\begin{equation}
    \Phi(\bar{\rho}_t, \kappa_t) = 1 + \kappa(\bar{\rho}_t) + \gamma(HHI_t),
    \label{eq:phi}
\end{equation}
where $\kappa(\bar{\rho}_t)$ is the correlation penalty and $\gamma(HHI_t)$
is the concentration penalty.

\subsubsection*{Correlation Penalty}

Let $\bar{\rho}_t$ denote the average pairwise Pearson correlation of daily
log-returns across all collateral positions, estimated over a 60-day rolling
window. The correlation penalty is:
\begin{equation}
    \kappa(\bar{\rho}_t) =
    \begin{cases}
        0 & \bar{\rho}_t \leq 0.4 \\[4pt]
        0.05 \cdot \dfrac{\bar{\rho}_t - 0.4}{0.6} & 0.4 < \bar{\rho}_t \leq 1
    \end{cases}
    \label{eq:kappa_layer2}
\end{equation}
This linearly ramps from 0\% to 5\% as average correlation rises from 0.4
to 1.0, reflecting diminishing diversification in stress.

\subsubsection*{Concentration Penalty}

Portfolio concentration is measured by the Herfindahl--Hirschman Index (HHI):
\begin{equation}
    HHI_t = \sum_{i=1}^{N} w_i^2,
    \label{eq:hhi}
\end{equation}
where $w_i$ are the adjusted value weights. The concentration penalty is:
\begin{equation}
    \gamma(HHI_t) =
    \begin{cases}
        0                                          & HHI_t \leq 0.10 \\[4pt]
        0.10 \cdot \dfrac{HHI_t - 0.10}{0.90}    & 0.10 < HHI_t \leq 1
    \end{cases}
    \label{eq:gamma}
\end{equation}
This ramps from 0\% to 10\% as HHI rises from a well-diversified portfolio
(0.10) to a fully concentrated single-name portfolio (1.0).

\subsection{Regime-Adjusted Portfolio Margin}
\label{subsec:regime_layer2}

Portfolio margin requirements are adjusted dynamically according to the
prevailing volatility regime. The regime classification (Normal, Elevated,
Stress) follows the three-state VSTOXX-based rule defined in
\texttt{trb\_method.tex} Section~5.2.

In each regime, a multiplicative regime scalar $R_t$ is applied to the
portfolio VaR:
\begin{equation}
    \MR_t^{(2),\mathrm{regime}} = R_t \cdot \MR_t^{(2)},
    \label{eq:mr_regime}
\end{equation}
with regime scalars:

\begin{table}[H]
\centering
\caption{Regime Scalars for Portfolio Margin Overlay (Layer~2)}
\label{tab:regime_scalars}
\begin{tabular}{lcc}
\toprule
\textbf{Regime} & \textbf{VSTOXX Level} & \textbf{Regime Scalar $R_t$} \\
\midrule
Normal   & $<$ 20  & 1.00 \\
Elevated & 20--30  & 1.15 \\
Stress   & $>$ 30  & 1.35 \\
\bottomrule
\end{tabular}
\end{table}

\noindent The asymmetric persistence rules for regime transitions (immediate
downgrade, 5-day confirmation for upgrade) defined in \texttt{trb\_method.tex}
Section~5.4 apply equally to the regime scalar.

\subsection{Single-Name Cap}

Regardless of the portfolio VaR result, no single issuer may contribute
more than 20\% of the adjusted collateral value in the portfolio overlay
computation. Positions exceeding this cap are split: the capped portion
is treated as a concentrated single-name exposure and assessed at the
Layer~1 rate without diversification benefit.

% ══════════════════════════════════════════════════════════════════════════════
\section{Layer Interaction and Effective Margin Requirement}
\label{sec:interaction}

\subsection{The Floor Constraint}

The two-layer architecture enforces the following ordering at all times:
\begin{equation}
    \MR_t = \max\!\left(\MR_t^{(1)},\; \MR_t^{(2),\mathrm{regime}}\right)
    \geq \MR_t^{(1)}.
    \label{eq:floor_constraint}
\end{equation}

This means:
\begin{itemize}[leftmargin=1.5em]
    \item In \textbf{Normal} regimes with a well-diversified portfolio
    ($\bar{\rho}_t$ low, $HHI_t$ low), the portfolio VaR may fall
    below the Layer~1 aggregate, but the effective requirement is still
    $\MR_t^{(1)}$ -- the regulatory floor is never breached.
    \item In \textbf{Stress} regimes or concentrated portfolios, the Layer~2
    overlay will exceed the Layer~1 floor and the portfolio-level figure
    drives the requirement.
\end{itemize}

\subsection{Decomposition of the Effective Margin Requirement}

For reporting and client communication purposes, the effective margin
requirement is decomposed as:

\begin{equation}
    \MR_t = \underbrace{\MR_t^{(1)}}_{\text{Reg T floor}}
    \;+\;
    \underbrace{\max\!\left(0,\; \MR_t^{(2),\mathrm{regime}} - \MR_t^{(1)}\right)}_{\text{Portfolio overlay add-on}}.
    \label{eq:mr_decomp}
\end{equation}

\noindent This decomposition is used in client-facing margin statements and
in the internal capital attribution to the lending desk.

\subsection{Link to the LTV Framework}

The margin requirement $\MR_t$ defined in equation~\eqref{eq:mr_total}
translates directly into the LTV threshold structure of \texttt{trb\_method.tex}
as follows. Let $E_t = V_t^{\mathrm{adj}} - L_t$ be the client equity. The
LTV-equivalent of the effective margin requirement is:
\begin{equation}
    \LTV^{\mathrm{warn}}_t
    = 1 - \frac{\MR_t}{V_t^{\mathrm{adj}}},
    \label{eq:ltv_link}
\end{equation}
i.e.\ the warning LTV threshold is set such that the client equity at
the threshold equals exactly the effective margin requirement. The
liquidation threshold is derived from this via an additional buffer
calibrated to the liquidation horizon, as specified in
\texttt{trb\_method.tex} Section~5.

% ══════════════════════════════════════════════════════════════════════════════
\section{Margin Call Process}
\label{sec:margin_call}

\subsection{Trigger Hierarchy}

Margin calls may be triggered at either layer:
\begin{enumerate}[leftmargin=1.5em]
    \item \textbf{Layer~1 trigger:} Client equity on any single position falls
    below the instrument-level maintenance margin $\MM_{i,t}$ (Section~\ref{subsec:maintenance_margin}).
    \item \textbf{Layer~2 trigger:} Portfolio-level LTV exceeds the
    regime-adjusted warning threshold $\LTV^{\mathrm{warn},\star}_t$ as
    defined in \texttt{trb\_method.tex} Section~6.
\end{enumerate}

\noindent A Layer~1 trigger requires cure of the breaching instrument
irrespective of portfolio-level LTV. A Layer~2 trigger requires
portfolio-level cure. Both triggers may be active simultaneously.

\subsection{Cure Mechanisms}

Three cure mechanisms are available (consistent with \texttt{trb\_method.tex}
Section~6.2):
\begin{itemize}[leftmargin=1.5em]
    \item \textbf{Cash deposit / partial loan repayment:} Reduces $L_t$,
    directly improving both layer computations.
    \item \textbf{Additional eligible collateral:} Increases $V_t^{\mathrm{adj}}$
    and reduces Layer~2 concentration (if diversifying) or Layer~1 shortfall.
    \item \textbf{Partial position reduction:} Reduces both exposure and,
    where the sold position is the binding concentration name, reduces
    $HHI_t$ and $\Phi$.
\end{itemize}

% ══════════════════════════════════════════════════════════════════════════════
\section{Governance and Validation}
\label{sec:governance}

\subsection{Review Schedule}

\begin{table}[H]
\centering
\caption{Governance Review Schedule}
\begin{tabular}{ll}
\toprule
\textbf{Activity} & \textbf{Frequency} \\
\midrule
Layer~1 rate schedule review                     & Quarterly \\
Layer~1 backtesting (Kupiec POF per asset class) & Quarterly \\
Layer~2 covariance model recalibration (EWMA)    & Monthly \\
Layer~2 regime scalar review                     & Semi-annual \\
Concentration / correlation penalty review       & Quarterly \\
Full methodology review                          & Annual \\
Lender margin report (Layer~1 + Layer~2 split)   & Daily \\
\bottomrule
\end{tabular}
\end{table}

\subsection{Backtesting}

Layer~1 haircut adequacy is validated using the Kupiec Proportion of
Failures (POF) test as specified in \texttt{trb\_method.tex} Section~9.1.
For Layer~2, the portfolio VaR is backtested by comparing the predicted
$\VaR_\alpha$ against realised daily portfolio P\&L. An exceedance rate
$\hat{p} = x/T$ significantly above $1-\alpha$ triggers model review.
Both the Kupiec and Christoffersen conditional coverage tests are applied.

\subsection{Model Independence and Override}

Layer~1 and Layer~2 are computed by independent model components. In the
event of a Layer~2 model failure or data unavailability, the effective
margin requirement falls back to $\MR_t^{(1)}$ automatically, ensuring
the Reg~T floor is always enforced. Manual overrides to either layer
require dual approval by Risk Management and the Head of Credit Risk,
with written justification retained for audit purposes.

% ══════════════════════════════════════════════════════════════════════════════
\appendix

\section{Notation Summary}
\label{app:notation}

\begin{longtable}{cl}
\toprule
\textbf{Symbol} & \textbf{Description} \\
\midrule
\endhead
$\MR_t$                          & Effective margin requirement at $t$ \\
$\MR_t^{(1)}$                    & Layer~1 (instrument-level) aggregate margin requirement \\
$\MR_t^{(2)}$                    & Layer~2 (portfolio-level) margin requirement \\
$\IM_{i,t}$                      & Initial margin requirement for instrument $i$ at $t$ \\
$\MM_{i,t}$                      & Maintenance margin requirement for instrument $i$ at $t$ \\
$m_i^{\IM}$                      & Initial margin rate for instrument $i$ \\
$m_i^{\MM}$                      & Maintenance margin rate for instrument $i$ \\
$P_{i,t}$                        & Price of instrument $i$ at $t$ \\
$Q_i$                            & Quantity held of instrument $i$ \\
$h_i$                            & Haircut for instrument $i$ \\
$w_i$                            & Adjusted value weight of instrument $i$ \\
$\hat{\Sigma}_t$                 & EWMA covariance matrix at $t$ \\
$\hat{\sigma}_{p,t}$             & Portfolio volatility at $t$ \\
$\bar{\rho}_t$                   & Average pairwise correlation at $t$ \\
$\kappa(\bar{\rho}_t)$           & Correlation penalty \\
$HHI_t$                          & Herfindahl--Hirschman Index at $t$ \\
$\gamma(HHI_t)$                  & Concentration penalty \\
$\Phi(\bar{\rho}_t, \kappa_t)$   & Combined concentration and correlation multiplier \\
$R_t$                            & Regime scalar \\
$\VaR_\alpha$                    & Value-at-Risk at confidence level $\alpha$ \\
$z_\alpha$                       & Standard normal quantile at level $\alpha$ \\
$\Delta t$                       & Liquidation horizon (years) \\
$\lambda$                        & EWMA decay factor \\
$\LTV^{\mathrm{warn}}_t$        & Warning LTV threshold at $t$ (linked to $\MR_t$) \\
$p^*$                            & Target loss probability tolerance \\
\bottomrule
\end{longtable}

\section{Worked Example: Two-Layer Margin Requirement}
\label{app:example}

\subsection*{Portfolio}
\begin{itemize}[leftmargin=1.5em]
    \item 1,000 shares large-cap equity: $P = \text{EUR}\;80$, $h = 20\%$,
    $\sigma = 18\%$ p.a., $m^{\IM} = 50\%$.
    \item EUR~20,000 nominal 3Y AA government bond: price 101\%,
    $h = 2\%$, $\sigma = 3\%$ p.a., $m^{\IM} = 5\%$.
    \item 500 shares mid-cap equity: $P = \text{EUR}\;40$, $h = 30\%$,
    $\sigma = 28\%$ p.a., $m^{\IM} = 60\%$.
\end{itemize}

\subsection*{Layer 1 Computation}
\[
\MR^{(1)} = 0.50 \times 80{,}000 + 0.05 \times 20{,}200 + 0.60 \times 20{,}000
= 40{,}000 + 1{,}010 + 12{,}000 = \text{EUR}\;53{,}010.
\]

\subsection*{Layer 2 Computation}
Portfolio adjusted value (from \texttt{trb\_method.tex} Appendix~B):
$V^{\mathrm{adj}} = \text{EUR}\;97{,}796$, $L = \text{EUR}\;55{,}000$.

Portfolio weights: $w_1 = 65.4\%$, $w_2 = 20.2\%$, $w_3 = 14.3\%$.

$\hat{\sigma}_p \approx 16.0\%$ p.a.\ (from \texttt{trb\_method.tex}).

\[
\VaR_{0.99} = 2.326 \times 0.160 \times \sqrt{5/252} \times 97{,}796
\approx \text{EUR}\;16{,}234.
\]

$HHI = 0.654^2 + 0.202^2 + 0.143^2 = 0.428 + 0.041 + 0.020 = 0.489$,
so $\gamma = 0.10 \times (0.489 - 0.10)/0.90 = 4.3\%$.

$\bar{\rho} = 0.30 < 0.40$, so $\kappa = 0\%$.

$\Phi = 1 + 0 + 0.043 = 1.043$.

Normal regime: $R = 1.00$.

\[
\MR^{(2)} = 1.043 \times 16{,}234 = \text{EUR}\;16{,}932.
\]

\subsection*{Effective Margin Requirement}
\[
\MR = \max(53{,}010,\; 16{,}932) = \textbf{EUR}\;53{,}010.
\]
Layer~1 is binding in this case (Normal regime, moderate portfolio).
The portfolio overlay add-on is zero. In a Stress regime with
$R = 1.35$ and elevated correlation, Layer~2 would compute
$\approx \text{EUR}\;22{,}900$ -- still below Layer~1, confirming
the Reg~T floor dominates for this relatively conservative portfolio.

\subsection*{Link to LTV Warning Threshold}
\[
\LTV^{\mathrm{warn}} = 1 - \frac{53{,}010}{97{,}796} = 45.8\%.
\]
This warning LTV is more conservative than the regime-based threshold
in \texttt{trb\_method.tex} (70\% in Normal), confirming the
instrument-level floor is the binding constraint here.

\end{document}
